% !TeX root = ../tesis.tex

\section{Conclusiones}

\section{Trabajo futuro}

\subsection{Sobre \CIP-133}
TODO

\subsection{Sobre `zk-bridge-operator`}

El comando \texttt{zk-bridge-operator tx prove <transaction-hash>} genera hoy un \textit{bundle} con dos pruebas distintas para el mismo \texttt{tx\_id} canónico de \Cardano. La primera, \texttt{snapshot\_membership}, sigue la ruta \textit{legacy} hoy expuesta por el \MithrilAggregator para \CardanoTransactions y prueba que la transacción pertenece al snapshot certificado. La segunda, \TxSetUpdateCircuit, produce la prueba de actualización del \SMT local del bridge.

En el estado actual de la rama \texttt{preview}, sin embargo, el operador no construye aún un \textit{witness} completamente general para la segunda prueba. La función \texttt{single\_insert\_empty\_tree\_witness} arma un caso canónico de inserción sobre árbol vacío: parte de la raíz del \SMT completamente vacío, usa el camino inducido por \texttt{tx\_id} y calcula la raíz resultante al cambiar la hoja final de $0$ a $1$. Esto es suficiente para la prueba de concepto actual y no altera la semántica matemática del circuito, pero claramente estrecha el conjunto de instancias que el operador sabe producir hoy.

La consecuencia conceptual es nítida. \TxSetUpdateCircuit ya prueba la relación general \TxSetUpdateRelation: cualquier par $(r_{\mathrm{in}}, r_{\mathrm{out}})$ compatible con un mismo camino y un mismo vector de hermanos entra en su lenguaje. Lo que permanece como trabajo futuro no es generalizar el circuito, sino generalizar la capa \Offchain del operador para que construya \textit{witnesses} de actualización sobre un \SMT arbitrario mantenido a lo largo de ejecuciones sucesivas, y no sólo sobre el caso canónico de primer inserto.

\subsection{Sobre `preview testnet` de \Cardano}